\subsection{DECOMPOSITIONS OF DECOMPOSABLE ALGEBRAS}

If g is a Lie subalgebra of \(\mathfrak{gl}(V)\) with radical r, the set of nilpotent elements of r is a nilpotent ideal of g, the largest nilpotency ideal of the identity representation of g (Chap. I, §5, no. 3, Cor. 6 of Th. 1). In this paragraph, we shall denote this ideal by \(\mathfrak{n}_{\mathrm{V}}(\mathfrak{g})\) . It contains the nilpotent radical \([g,g]\cap r\) of g (Chap. I, §5, no. 3, Th. 1).

PROPOSITION 5. Let g be a decomposable nilpotent Lie subalgebra of \(\mathfrak{gl}(V)\) . Let t be the set of semi-simple elements of g. Then t is a central subalgebra of g, and g is the product of t and \(\mathfrak{n}_{\mathrm{V}}(\mathfrak{g})\) as Lie algebras.

If \(x \in t\) , \(ad_{g}x\) is semi-simple and nilpotent, hence zero, so that x is central in g. Consequently, t is an ideal of g, and \(t \cap \mathfrak{n}_{\mathrm{V}}(\mathfrak{g}) = 0\) . Since g is decomposable, \(\mathfrak{g} = \mathfrak{t} + \mathfrak{n}_{\mathrm{V}}(\mathfrak{g})\) , hence the proposition.

PROPOSITION 6. Let g be a decomposable Lie subalgebra of \(\mathfrak{gl}(\mathrm{V})\) . Let T be the set of commutative subalgebras of g consisting of semi-simple elements, and \(T_{1}\) the set of maximal elements of T. Let H be the set of Cartan subalgebras of g.

\begin{enumerate}
\def\labelenumi{(\roman{enumi})}
\item
  For \(\mathfrak{h} \in \mathcal{H}\) , let \(\varphi(\mathfrak{h})\) be the set of semi-simple elements of \(\mathfrak{h}\) . Then \(\varphi(\mathfrak{h}) \in \mathcal{T}_1\) .
\item
  For \(\mathfrak{t} \in \mathcal{T}_1\) , let \(\psi(\mathfrak{t})\) be the commutant of \(\mathfrak{t}\) in \(\mathfrak{g}\) . Then \(\psi(\mathfrak{t}) \in \mathscr{H}\) .
\item
  The maps \(\varphi\) and \(\psi\) are inverse bijections from \(\mathcal{H}\) to \(\mathcal{T}_1\) and from \(\mathcal{T}_1\) to \(\mathcal{H}\) .
\item
  If \(k\) is algebraically closed, \(\mathrm{Aut}_e(\mathfrak{g})\) operates transitively on \(\mathcal{T}_1\) .
\end{enumerate}

Let \(\mathfrak{h} \in \mathcal{H}\) , and put \(\mathfrak{t} = \varphi(\mathfrak{h})\) . By Prop. 5 and Cor. 2 of Prop. 3, \(\mathfrak{t} \in \mathcal{T}\) and \(\mathfrak{h} = \mathfrak{t} \times \mathfrak{n}_{\mathrm{V}}(\mathfrak{h})\) . For any subalgebra \(\mathfrak{u}\) of \(\mathfrak{g}\) , we denote by \(\psi(\mathfrak{u})\) the commutant of \(\mathfrak{u}\) in \(\mathfrak{g}\) . Then \(\mathfrak{h} \subset \psi(\mathfrak{t})\) , and \(\psi(\mathfrak{t}) \subset \mathfrak{g}^0(\mathfrak{h})\) since the elements of \(\mathfrak{n}_{\mathrm{V}}(\mathfrak{h})\) are nilpotent, so \(\mathfrak{h} = \psi(\mathfrak{t})\) . If \(\mathfrak{t}' \in \mathcal{T}\) and \(\mathfrak{t} \subset \mathfrak{t}'\) , we have \(\mathfrak{t}' \subset \psi(\mathfrak{t}) = \mathfrak{h}\) so \(\mathfrak{t}' = \mathfrak{t}\) , and hence \(\mathfrak{t} \in \mathcal{T}_1\) .

Let \(\mathfrak{t} \in \mathcal{T}_1\) , and put \(\mathfrak{c} = \psi(\mathfrak{t})\) . Let \(\mathfrak{h}\) be a Cartan subalgebra of \(\mathfrak{c}\) . By §2, no. 3, Prop. 10, \(\mathfrak{h} \in \mathcal{H}\) and \(\mathfrak{t} \subset \mathfrak{h}\) . Put \(\mathfrak{t}_1 = \varphi(\mathfrak{h}) \in \mathcal{T}\) . Then \(\mathfrak{t} \subset \mathfrak{t}_1\) so \(\mathfrak{t} = \mathfrak{t}_1\) , and \(\mathfrak{h} = \psi(\mathfrak{t}_1) = \psi(\mathfrak{t}) = \mathfrak{c}\) by the preceding. Thus, \(\psi(\mathfrak{t}) \in \mathcal{H}\) , and \(\varphi(\psi(\mathfrak{t})) = \mathfrak{t}\) .

We have thus proved (i), (ii) and (iii). Assume that \(k\) is algebraically closed. Since \(\mathrm{Aut}_e(\mathfrak{g})\) operates transitively on \(\mathcal{H}(\S 3, \text{no. } 2, \text{Th. } 1)\) , \(\mathrm{Aut}_e(\mathfrak{g})\) operates transitively on \(\mathcal{T}_1\) .

COROLLARY 1. The Cartan subalgebras of \(\mathfrak{g}\) are the centralizers of the regular semi-simple elements of \(\mathfrak{g}\) .

If \(x \in \mathfrak{g}\) is regular, \(\mathfrak{g}^0(x)\) is a Cartan subalgebra of \(\mathfrak{g}\) (§2, no. 3, Th. 1 (i)); moreover, if \(x\) is semi-simple \(\mathfrak{g}^0(x)\) is the centralizer of \(x\) in \(\mathfrak{g}\) . Conversely, let \(\mathfrak{h}\) be a Cartan subalgebra of \(\mathfrak{g}\) . There exists \(\mathfrak{t} \in \mathcal{T}_1\) such that \(\mathfrak{h} = \psi(\mathfrak{t})\) . By §1, no. 2, Prop. 7, there exists \(x \in \mathfrak{t}\) such that \(\mathfrak{h} = \mathfrak{g}^0(x)\) ; since \(x \in \mathfrak{t}\) , \(\mathfrak{g}^0(x) = \mathfrak{g}_0(x)\) . By §3, no. 3, Th. 2 (ii), \(x\) is regular.

COROLLARY 2. Assume in addition that g is solvable. Then:

\begin{enumerate}
\def\labelenumi{(\roman{enumi})}
\item
  The subgroup of \(\operatorname{Aut}(\mathfrak{g})\) consisting of the \(e^{\mathrm{ad}x}, x \in \mathcal{C}^{\infty}\mathfrak{g}\) (cf.~§3, no. 4), operates transitively on \(\mathcal{T}_1\) .
\item
  If \(\mathfrak{t} \in \mathcal{T}_1\) , \(\mathfrak{g}\) is the semi-direct product of \(\mathfrak{t}\) and \(\mathfrak{n}_{\mathrm{V}}(\mathfrak{g})\) .
\end{enumerate}

Assertion (i) follows from the fact that the group of the \(e^{\mathrm{ad}x}\) , \(x \in \mathcal{C}^{\infty}\mathfrak{g}\) , operates transitively on \(\mathcal{H}(\S 3\) , no. 4, Th. 3).

We prove (ii). Let \(t \in T_{1}\) , and let \(\mathfrak{h} = \psi(t)\) be the corresponding Cartan subalgebra of g. In view of Prop. 5, \(\mathfrak{h} = \mathfrak{t} + \mathfrak{n}_{\mathrm{V}}(\mathfrak{h}) \subset \mathfrak{t} + \mathfrak{n}_{\mathrm{V}}(\mathfrak{g})\) . On the other hand, \(g = h + [g, g]\) (§2, no. 1, Cor. 3 of Prop. 4) and \([\mathfrak{g}, \mathfrak{g}] \subset \mathfrak{n}_{\mathrm{V}}(\mathfrak{g})\) , so \(\mathfrak{g} = \mathfrak{t} + \mathfrak{n}_{\mathrm{V}}(\mathfrak{g})\) . But it is clear that \(\mathfrak{t} \cap \mathfrak{n}_{\mathrm{V}}(\mathfrak{g}) = \{0\}\) . The algebra g is thus the semi-direct product of t and the ideal \(\mathfrak{n}_{\mathrm{V}}(\mathfrak{g})\) .

PROPOSITION 7. Let \(\mathfrak{g}\) be a decomposable Lie subalgebra of \(\mathfrak{gl}(\mathrm{V})\) .

\begin{enumerate}
\def\labelenumi{(\roman{enumi})}
\item
  There exists a Lie subalgebra \(\mathfrak{m}\) of \(\mathfrak{g}\) , reductive in \(\mathfrak{gl}(\mathrm{V})\) , such that \(\mathfrak{g}\) is the semi-direct product of \(\mathfrak{m}\) and \(\mathfrak{n}_{\mathrm{V}}(\mathfrak{g})\) .
\item
  Any two Lie subalgebras of \(\mathfrak{g}\) with the properties in (i) are conjugate under \(\mathrm{Aut}_e(\mathfrak{g})\) .
\end{enumerate}

The radical r of g is decomposable (no. 2, Cor. 2 of Prop. 4). By Cor. 2 of Prop. 6, there exists a commutative subalgebra t of r, consisting of semi-simple elements, such that \(\mathfrak{r} = \mathfrak{t} \oplus \mathfrak{n}_{\mathrm{V}}(\mathfrak{r})\) . Since \(ad_{g}t\) consists of semi-simple elements, g is the direct sum of \([t, g]\) and the centralizer z of t (Chap. I, §3, no. 5, Prop. 6). Since \([t, g] \subset r\) , \(g = z + r\) . Consequently, if s is a Levi subalgebra of z (Chap. I, §6, no. 8), \(g = s + r\) , so s is a Levi subalgebra of g. Put \(m = s \oplus t\) . Since \([s, t] = \{0\}\) , m is a Lie subalgebra of g, reductive in \(\mathfrak{gl}(V)\) by Chap. I, §6, no. 5, Th. 4. Moreover,

\[
\mathfrak {g} = \mathfrak {s} \oplus \mathfrak {r} = \mathfrak {s} \oplus \mathfrak {t} \oplus \mathfrak {n} _ {\mathrm{V}} (\mathfrak {r}) = \mathfrak {s} \oplus \mathfrak {t} \oplus \mathfrak {n} _ {\mathrm{V}} (\mathfrak {g}) = \mathfrak {m} \oplus \mathfrak {n} _ {\mathrm{V}} (\mathfrak {g})
\]

since \(\mathfrak{n}_{\mathrm{V}}(\mathfrak{g}) = \mathfrak{n}_{\mathrm{V}}(\mathfrak{r})\) . Hence (i).

Now let \(\mathfrak{m}'\) be a Lie subalgebra of \(\mathfrak{g}\) complementary to \(\mathfrak{n}_{\mathrm{V}}(\mathfrak{g})\) and reductive in \(\mathfrak{gl}(\mathrm{V})\) . We show that \(\mathfrak{m}'\) is conjugate to \(\mathfrak{m}\) under \(\operatorname{Aut}_e(\mathfrak{g})\) . We have \(\mathfrak{m}' = \mathfrak{s}' \oplus \mathfrak{t}'\) , where \(\mathfrak{s}' = [\mathfrak{m}', \mathfrak{m}']\) is semi-simple and the centre \(\mathfrak{t}'\) of \(\mathfrak{m}'\) consists of semi-simple elements. Then \(\mathfrak{r} = \mathfrak{t} \oplus \mathfrak{n}_{\mathrm{V}}(\mathfrak{g}) = \mathfrak{t}' \oplus \mathfrak{n}_{\mathrm{V}}(\mathfrak{g})\) . In view of Cor. 2 of Prop. 6, we are reduced to the case \(\mathfrak{t} = \mathfrak{t}'\) . Then \(\mathfrak{s}' \subset \mathfrak{z}\) ; since \(\dim \mathfrak{s}' = \dim \mathfrak{s}\) , \(\mathfrak{s}'\) is a Levi subalgebra of \(\mathfrak{z}\) . By Chap. I, §6, no. 8, Th. 5, there exists \(x \in \mathfrak{n}_{\mathrm{V}}(\mathfrak{z})\) such that \(e^{\operatorname{ad} x}(\mathfrak{s}) = \mathfrak{s}'\) ; since \(x\) commutes with \(\mathfrak{t}\) , we also have \(e^{\operatorname{ad} x}(\mathfrak{t}) = \mathfrak{t}\) .

